\section{Holomorphic coefficient uniqueness on the original frame domain}
\label{sec:holomorphic-uniqueness}

This chapter reconstructs the negative-frame uniqueness argument used by
the original theta extension, together with its representation-theoretic
consumers.  All source identifications refer to commit
\texttt{7cbf7a0c32fca1badf05931879a37c3f8a2b03b4}.
The arguments are on the supplied original objects.  An existential
auxiliary frame factor is constructed where the source asks for one;
the original Hermitian forms, discrete-series models and coefficient
maps are not changed.

\subsection*{The exact finite-dimensional dictionary}

Let \(I\) be a finite set.  At each \(\tau\in I\) fix an \(n\)-by-\(n\)
complex matrix \(H_\tau\), an integer \(0\le q_\tau\le n\), and a matrix
\(x_{0,\tau}\) satisfying the source's \(\mathrm{OrthFr}\) condition
\[
 x_{0,\tau}^{*}H_\tau x_{0,\tau}=J_\tau,\qquad
 J_\tau=-S_\tau,\qquad
 S_\tau=\operatorname{diag}(s_{\tau,a}),\quad
 s_{\tau,a}=
 \begin{cases}1&a<q_\tau,\\-1&a\ge q_\tau.\end{cases}
 \tag{HCU.1}
\]
Indices start at zero.  For a matrix \(A\), let \(\Delta_k(A)\) be the
determinant of its leading \(k\)-by-\(k\) submatrix, with
\(\Delta_0(A)=1\), and put
\(s_{\tau,[k]}=\prod_{a<k}s_{\tau,a}\).
The actual source domain is
\[
 \mathcal D_\tau=\{x:\det x\ne0,\quad
 s_{\tau,[k]}\operatorname{Re}\Delta_k(-x^*H_\tau x)>0
 \ (0\le k\le n)\},\qquad
 \mathcal D=\prod_{\tau\in I}\mathcal D_\tau .
 \tag{HCU.2}
\]
Let \(B\) be the product of the invertible upper-triangular groups,
\(G=\prod_\tau U(H_\tau)\), and
\[
 \chi_w(b)=\prod_{\tau,a}b_{\tau,aa}^{\,w_{\tau,a}},
 \quad w_{\tau,a}\in\mathbb Z,\qquad
 D(s)_\tau=\operatorname{diag}(e^{is_{\tau,a}}),\quad
 t(s)=x_0D(s)x_0^{-1}.
 \tag{HCU.3}
\]
Here and below \(s\) in \(D(s)\) is a tuple of real angles, distinct from
the fixed signs in (HCU.1).  Write
\(\chi_w(s)=\prod_{\tau,a}e^{iw_{\tau,a}s_{\tau,a}}\).
The source's \(\mathrm{HolFn}(H,q,x_0,w,f)\) means exactly that there
is a holomorphic function \(P\) on \(\mathcal D\) such that
\[
 P(xb)=\chi_w(b)P(x),\qquad f(u)=P(ux_0)\quad(u\in G).
 \tag{HCU.4}
\]
Holomorphy means complex differentiability on this open subset of the
finite-dimensional matrix space.  No extension to all matrices is
required.  Empty products and zero-dimensional matrix spaces are allowed.

% BEGIN HCU1
\paragraph{HCU1: a smooth triangular factor on the original domain.}
\label{claim:HCU1}
The matrices \(H_\tau,x_{0,\tau}\) are nonsingular, and \(H_\tau\) is
Hermitian.  There is a real-smooth map
\(\beta_\tau:\mathcal D_\tau\to B_n\) such that
\[
 u_\tau(x)=x\beta_\tau(x)^{-1}x_{0,\tau}^{-1}\in U(H_\tau),
 \qquad x=u_\tau(x)x_{0,\tau}\beta_\tau(x).
 \tag{HCU.5}
\]
This is a witness for the literal
\texttt{smooth\_gramSchmidt} existential statement.

\emph{Proof.}
Taking determinants in (HCU.1) proves nonsingularity, and solving that
equation for \(H_\tau\) proves Hermitian symmetry.  Omit \(\tau\).
For \(A=-x^*Hx\), all leading minors are nonzero.  Gaussian elimination
without row exchanges gives
\[
 A=L\operatorname{diag}(d_0,\ldots,d_{n-1})L^*,
 \quad L\text{ lower unitriangular},\qquad
 d_a=\frac{\Delta_{a+1}(A)}{\Delta_a(A)}\in\mathbb R^\times .
 \tag{HCU.6}
\]
For completeness, isolate the nonzero first pivot \(d=A_{00}\).
Subtracting the outer product of the first column and row, divided by
\(d\), leaves its Hermitian Schur complement.  Its leading \(k\)-minor
is \(\Delta_{k+1}(A)/d\).  Induction constructs (HCU.6), and comparison
of leading determinants gives its displayed pivot formula.  Every
entry of \(L\) and every \(d_a\) is a rational expression in entries of
\(A\), with the already nonzero pivots as denominators.

The signs in (HCU.2) give \(s_a d_a>0\).  Set
\[
 \beta(x)=\operatorname{diag}(\sqrt{|d_a|})\,L^* .
 \tag{HCU.7}
\]
Then \(\beta\) is upper triangular with positive diagonal, and
\(A=\beta^*S\beta\).  Consequently
\[
 (x\beta^{-1})^*H(x\beta^{-1})=-S=x_0^*Hx_0,
\]
which proves (HCU.5).  The rational formulas, complex conjugation and square root on
the positive real axis show real smoothness on the entire domain.
For \(n=0\) all matrices are the unique empty matrix and all formulas
have their empty-product meaning.  The source requires an existential
factor with (HCU.5); it does not prescribe the particular diagonal
sign convention of that factor.  Thus (HCU.7) changes no stored datum.
% END HCU1

% BEGIN HCU2
\paragraph{HCU2: openness, stability and connectedness of the actual domain.}
\label{claim:HCU2}
The set \(\mathcal D\) is open, nonempty and path connected.  It contains
\(x_0\) and is stable under \(x\mapsto uxb\), \(u\in G,b\in B\).
In particular it has exactly the properties used in
\texttt{frDoms\_open\_preconnected\_stable}.

\emph{Proof.}
The finitely many determinant inequalities in (HCU.2) are open
conditions.  For upper-triangular \(b\), its first \(k\) columns involve
only the first \(k\) rows, so
\[
 \Delta_k\bigl(b^*(-x^*Hx)b\bigr)
   =|\det b_{[k]}|^2\Delta_k(-x^*Hx).
 \tag{HCU.8}
\]
The multiplying scalar is positive.  Left multiplication by \(U(H)\)
does not change the Gram matrix.  These observations prove stability;
(HCU.1) proves \(x_0\in\mathcal D\).  HCU1 gives the surjection
\((u,b)\mapsto ux_0b\) from \(G\times B\) to \(\mathcal D\).

Here is a direct path-connectedness proof for its two factors.
Conjugation by \(x_0\) identifies \(U(H)\) with \(U(J)\).
For \(g\in U(J)\), put \(A=g^*g\), \(p=A^{1/2}\), \(k=gp^{-1}\).
The positive square root is defined by unitary diagonalization of the
positive Hermitian matrix \(A\).  From \(g^{-1}=Jg^*J\) one obtains
\(JAJ=A^{-1}\), hence \(JpJ=p^{-1}\).  Thus \(k\) is both ordinarily
unitary and \(J\)-unitary, and therefore commutes with \(J\).
Choose a skew-Hermitian logarithm \(Y\) of \(k\), separately on its
two \(J\)-eigenspaces, and let \(X=\log p\), defined by the same
positive spectral calculus.  Then
\[
 JXJ=-X,\qquad
 [Y,J]=0,\qquad
 [0,1]\ni t\longmapsto e^{tY}e^{tX}\in U(J)
 \tag{HCU.9}
\]
joins \(1\) to \(g\).  These assertions follow on each eigenspace:
positive eigenvalues have real logarithms, unit eigenvalues have
imaginary logarithms, and conjugation by \(J\) takes the positive
spectral calculus of \(A\) to that of \(A^{-1}\).  Only the existence
of a path for each fixed \(g\) is used; no continuous global choice of
logarithms is being asserted.

An upper-triangular invertible matrix has the form \(d(1+N)\), with
\(d\) diagonal invertible and \(N\) strictly upper triangular.
For fixed logarithms \(\ell_a\) of its diagonal entries,
\(\operatorname{diag}(e^{t\ell_a})(1+tN)\) is a path from \(1\) to it.
Finite products and continuous images preserve path connectedness.
This proves the claim, also when a signature block, \(I\), or \(n\)
is empty.
% END HCU2

% BEGIN HCU3
\paragraph{HCU3: the local zero lemma with its exact hypotheses.}
\label{claim:HCU3}
Let \(\Omega\) be a neighborhood of the tuple \(1\) in
\(\prod_I M_n(\mathbb C)\), and let \(Q\) be a function on that
ambient space which is complex differentiable on \(\Omega\).
Suppose:
\begin{enumerate}
\item if \(y\in\Omega\), \(Q(y)=0\) and \(b\in B\), then \(Q(yb)=0\);
\item whenever \(y,y'\in\Omega\) and \(y'D(s)=D(s)y\), one has
\(Q(y')=Q(y)\);
\item \(Q(1)=0\).
\end{enumerate}
Then \(Q\) vanishes on a neighborhood of \(1\).  Only its restriction
to a small open ball inside \(\Omega\) is used analytically;
the first hypothesis is the literal ambient zero-set condition.

\emph{Proof.}
Use the maximum-entry norm and choose \(r>0\) with
\(B(1,r)\subset\Omega\).  For a tuple \(z\) put
\[
 L_\zeta(z)_{\tau,ac}=
 \begin{cases}\zeta^{a-c}z_{\tau,ac}&c<a,\\0&c\ge a,\end{cases}
 \qquad
 U(z)_{\tau,ac}=
 \begin{cases}z_{\tau,ac}&a\le c,\\0&a>c.\end{cases}
 \tag{HCU.10}
\]
If \(\|z\|<r\) and \(|\zeta|\le1\), then
\(\|L_\zeta(z)\|\le\|z\|\).
Thus \(h(\zeta)=Q(1+L_\zeta(z))\) is holomorphic on a neighborhood of
the closed unit disc.  For \(\zeta=e^{it}\),
\(D_\zeta=\operatorname{diag}(\zeta^a)\) gives
\[
 (1+L_\zeta(z))D_\zeta=D_\zeta(1+L_1(z)).
 \tag{HCU.11}
\]
Hypothesis 2 makes \(h\) constant on the unit circle.
The maximum-modulus principle applied to \(h-h(1)\), or the Cauchy
integral formula, makes it constant throughout the disc.
Since all exponents \(a-c\) in the lower part are positive,
\(L_0(z)=0\), and \(h(1)=h(0)=Q(1)=0\).

Now the polynomial map
\[
 \mathcal L(z)=(1+L_1(z))(1+U(z))
       =1+z+L_1(z)U(z)
 \tag{HCU.12}
\]
has derivative the identity at zero.  The finite-dimensional inverse
function theorem therefore maps a sufficiently small neighborhood of
zero onto a neighborhood of \(1\).  Shrink it to \(\|z\|<\min(r,1)\).
The diagonal entries \(1+z_{\tau,aa}\) are nonzero, so \(1+U(z)\in B\).
The first part of the proof and hypothesis 1 give
\(Q(\mathcal L(z))=0\) throughout that neighborhood.  This proves
exactly \texttt{Uniq.local\_zero}.  In a zero-dimensional ambient space,
the conclusion is already \(Q(1)=0\).
% END HCU3

% BEGIN HCU4
\paragraph{HCU4: global uniqueness with the original torus law.}
\label{claim:HCU4}
Assume \(\mathrm{HolFn}(H,q,x_0,w,f)\), together with the separate
left torus law
\[
 f(t(s)u)=\chi_w(s)f(u)\quad(u\in G),\qquad f(1)=0.
 \tag{HCU.13}
\]
Then \(f=0\) on \(G\).  This is the full original
\texttt{Uniq.eq\_zero\_of\_torus\_law} contract.

\emph{Proof.}
Let \(P\) be its section in (HCU.4).  For \(x=ux_0b\), HCU1 and
(HCU.13) give
\[
 P(t(s)x)=\chi_w(b)f(t(s)u)=\chi_w(s)P(x).
 \tag{HCU.14}
\]
Consider \(Q(y)=P(x_0y)\) on
\(\Omega=\{y:x_0y\in\mathcal D\}\), an open neighborhood of \(1\).
Its zero set is stable under right multiplication by \(B\).
If \(y'=D(s)yD(s)^{-1}\), then
\[
 Q(y')=P(t(s)x_0yD(s)^{-1})
      =\chi_w(s)\chi_w(D(s)^{-1})Q(y)=Q(y).
 \tag{HCU.15}
\]
The integer powers are well defined and the two characters cancel
exactly.  Also \(Q(1)=f(1)=0\).  HCU3 makes \(P\) vanish on an open
neighborhood of \(x_0\).

A holomorphic function on a connected open set that vanishes on a
nonempty open subset vanishes everywhere.  To see the needed
several-variable statement directly, cover a path by overlapping balls.
On a ball, restrict to a complex line through a point of an open zero
set; the one-variable power-series identity theorem gives vanishing
along that line wherever the line stays in the ball.  It follows that
the whole ball is zero, and finitely many overlapping balls propagate
the conclusion along the path.  HCU2 supplies such paths inside the
actual domain.  Therefore \(P=0\) on \(\mathcal D\), whence \(f=0\)
on \(G\).
% END HCU4

% BEGIN HCU5
\paragraph{HCU5: the full smooth Borel--holomorphic dictionary.}
\label{claim:HCU5}
Let \(\mathfrak g_\tau=\{X:X^*H_\tau+H_\tau X=0\}\).
Use the source right derivatives \(R_X\) and their complexification
\[
 c_H(Z)=-H^{-1}Z^*H,\qquad
 R_Z^{\mathbb C}
 =R_{(Z+c_H(Z))/2}
    +iR_{(Z-c_H(Z))/(2i)}.
 \tag{HCU.16}
\]
A function \(f\) on \(G\) has \(\mathrm{HolFn}(H,q,x_0,w,f)\) if and
only if it has the source \(\mathrm{SmFn}\) smoothness and
\[
 R^{\mathbb C}_{x_{0,\tau}E_{ac}x_{0,\tau}^{-1}}f(u)
   =\begin{cases}w_{\tau,a}f(u)&a=c,\\0&a<c\end{cases}
 \quad(a\le c,\ u\in G).
 \tag{HCU.17}
\]
Thus both \texttt{holFn\_of\_killed} and
\texttt{smFn\_and\_killedFn\_of\_holFn} hold on the literal source domain.

\emph{Proof.}
First, the source smoothness condition, which tests finite products of
real one-parameter exponentials at every \(u\), is ordinary smoothness
on \(G\).  Here is the chart verification.  On the \(H\)-skew matrices,
the Cayley map
\((1+X/2)(1-X/2)^{-1}\) is a smooth chart into \(U(H)\) near \(1\);
its inverse is \(2(g-1)(g+1)^{-1}\), and its derivative at zero is the
identity.  Given a real basis \(X_1,\ldots,X_d\) of \(\mathfrak g\),
the product \(\prod_j e^{t_jX_j}\) has the same derivative in this chart.
The inverse function theorem makes these product exponentials local
coordinates.  The source tests their left translates, so it supplies
all smooth coordinate restrictions.  Conversely a smooth group function
satisfies every such test.

The right differential of a smooth function is real-linear in \(X\).
The involution \(c_H\) has fixed space \(\mathfrak g\), and every
complex matrix is uniquely \(X+iY\) with \(X,Y\in\mathfrak g\).
Consequently (HCU.16) is its complex-linear extension.
For a fixed real angle tuple \(s\), equation (HCU.17) on diagonal
units gives the scalar ordinary differential equation along
\(u\,t(ts)\), hence
\[
 f(u\,t(s))=\chi_w(s)f(u).
 \tag{HCU.18}
\]
The integer weights ensure that this is the prescribed character of
the diagonal torus.

Use HCU1 to define
\(P(x)=\chi_w(\beta(x))f(u(x))\).
This is real smooth.  It has the stronger formula
\[
 P(ux_0b)=\chi_w(b)f(u)\qquad(u\in G,\ b\in B).
 \tag{HCU.19}
\]
Indeed, two such decompositions differ by an element of
\(U(J_\tau)\cap B_n\).  An upper-triangular \(J\)-unitary matrix is
diagonal with entries of modulus one: its first column has only its
first entry, its norm determines that entry's modulus, and its
orthogonality to the other columns kills the remainder of the first
row; induction finishes.  Thus the two factors differ by \(D(s)\),
and (HCU.18) makes their two expressions in (HCU.19) identical.

It remains to prove complex differentiability; real smoothness alone
is insufficient.  At \(x=ux_0b\), write a tangent as
\(\dot x=ux_0Zb\).  Every complex matrix \(Z\), considered as a real
vector, is \(X+B_0\), where \(X\in\mathfrak u(J)\) and \(B_0\)
is upper triangular.  An explicit choice sets
\(X_{ac}=Z_{ac}\) for \(a>c\),
\(X_{ca}=-J_{aa}J_{cc}\overline{Z_{ac}}\), and \(X_{aa}=0\);
then \(B_0=Z-X\) is upper triangular.  Differentiating the factorized
curve
\(u e^{t x_0Xx_0^{-1}}x_0 e^{tB_0}b\) in (HCU.19) gives
\[
 dP_x(ux_0Zb)=\chi_w(b)
 \left(R_{x_0Xx_0^{-1}}f(u)
            +\sum_{\tau,a}w_{\tau,a}(B_0)_{\tau,aa}f(u)\right)
 =\chi_w(b)R^{\mathbb C}_{x_0Zx_0^{-1}}f(u).
 \tag{HCU.20}
\]
The last equality uses (HCU.17) for the whole upper-triangular
\(B_0\), by complex linearity.  Its right side is complex-linear in
\(Z\); multiplication by the fixed invertible frames identifies \(Z\)
complex-linearly with \(\dot x\).  Therefore the real derivative of
\(P\) is complex-linear at every point.  This is complex
differentiability, and proves the forward construction of (HCU.4).

Conversely let \(P\) satisfy (HCU.4).
The restrictions \(f(u)=P(ux_0)\) are smooth.
At \(ux_0\), its complex derivative in the tangent
\(ux_0E_{ac}\) is obtained from the curve \(ux_0e^{zE_{ac}}\).
For \(a<c\), its upper-triangular character is identically one; for
\(a=c\), it is \(e^{zw_{\tau,a}}\).
Differentiating gives (HCU.17), with precisely the complexification
(HCU.16).  This also proves (HCU.18) directly from (HCU.4).
% END HCU5

% BEGIN HCU6
\paragraph{HCU6: bounded integration preserves the original holomorphic type.}
\label{claim:HCU6}
Let \((\Omega,\nu)\) be a finite measure space, and suppose
\(f_o:G\to\mathbb C\) has the same \(\mathrm{HolFn}(H,q,x_0,w)\)
for every \(o\).  Suppose
\(|f_o(u)|\le C\) for all \(o,u\), and \(o\mapsto f_o(u)\) is
almost-everywhere strongly measurable for every fixed \(u\).
Then
\[
 u\longmapsto\int_\Omega f_o(u)\,d\nu(o)
 \tag{HCU.21}
\]
has the same \(\mathrm{HolFn}\).  This is the full
\texttt{holFn\_integral} statement.

\emph{Proof.}
Choose any section \(P_o\) for \(f_o\).  Its right transformation law
and HCU1 force the pointwise identity
\[
 P_o(x)=\chi_w(\beta(x))f_o(u(x)).
 \tag{HCU.22}
\]
In particular every section value is measurable in \(o\), regardless
of how the sections were chosen.  Near each \(x\in\mathcal D\),
the continuous factor \(\chi_w(\beta(y))\) is bounded: the diagonal
entries never vanish, including when some weights are negative.
Thus on a sufficiently small closed polydisc in \(\mathcal D\)
all \(P_o\) are bounded by a common finite constant.  Its constant
function on \(\Omega\) is integrable.

For clarity, holomorphy under this integral follows directly from
Cauchy's formula on a slightly larger polydisc.  Its coefficient
bounds give a uniformly summable majorant for the power series on
each smaller polydisc.  Integrate each coefficient; domination
justifies exchange of integral and series and gives a convergent
holomorphic power series for \(\int P_o\).
Measurability on the boundary products used in Cauchy's formula
follows by approximating their continuous parameter dependence with
finite grids: separate measurability in \(o\) and continuity in the
finite-dimensional compact parameter give measurable step
approximations with pointwise limits.  Hence Fubini is applicable.
The same argument with \(\max(C,0)\) covers the zero-measure or empty
index cases.  The right character law passes through the integral,
and evaluation at \(ux_0\) gives (HCU.21), proving the assertion.
% END HCU6

% BEGIN HCU7
\paragraph{HCU7: proportionality and a cyclic coefficient test.}
\label{claim:HCU7}
Two functions of the same \(\mathrm{HolFn}\) type, both satisfying
the left torus law in (HCU.13), and with values \(c,1\) at \(1\),
satisfy \(f=cg\) on \(G\).

For a single place, let \(\sigma\) be a strongly continuous unitary
irreducible representation of \(U(H)\) on a complex Hilbert space.
Let \(v\ne0\) be smooth and satisfy the original Borel equations
\[
 d\sigma^{\mathbb C}(x_0E_{ac}x_0^{-1})v
   =\delta_{ac}w_av\quad(a\le c).
 \tag{HCU.23}
\]
If \(l\) has torus character \(\chi_w\) and \(\langle l,v\rangle=0\),
then \(l=0\).  This closes the entire
\texttt{HolLowOnce.}\allowbreak\texttt{eq\_zero\_of\_orth} implication.

\emph{Proof.}
The difference \(f-cg\) satisfies HCU4; this proves proportionality.
For the second assertion, use the convention that the Hilbert inner
product is conjugate-linear in its first argument.  The smooth
coefficient \(F(u)=\langle l,\sigma(u)v\rangle\) satisfies (HCU.17)
by differentiating in the second argument, so HCU5 supplies its
literal \(\mathrm{HolFn}\).  Since \(\sigma(t(s))l=\chi_w(s)l\)
and \(|\chi_w(s)|=1\), unitarity gives
\[
 \langle l,\sigma(t(s)u)v\rangle
       =\chi_w(s)\langle l,\sigma(u)v\rangle .
 \tag{HCU.24}
\]
HCU4 makes this coefficient zero.  The closed span of
\(\{\sigma(u)v:u\in U(H)\}\) is a nonzero closed invariant subspace,
hence the whole Hilbert space.  Therefore \(l\) is orthogonal to
the whole space and is zero.  No assertion that all lowest compact
types were already irreducible was used in this test.
% END HCU7

\subsection*{The full original lowest compact type}

The next arguments concern the actual representation, not a replacement
model.  For a sign pattern \(\epsilon\), let
\(J_\epsilon=\operatorname{diag}(1_{\epsilon(a)}-1_{\neg\epsilon(a)})\),
\(K_\epsilon=U(n)\cap U(J_\epsilon)\), and write \(E_{ac}\) for matrix
units.  The source's \(\mathrm{IsKTypeOf}(\epsilon,\sigma,\Lambda,S)\)
says that \(S\) is finite dimensional and \(K_\epsilon\)-stable,
irreducible under this compact group, contains a nonzero torus vector
of weight \(\Lambda\), and every weight \(\mu\) in \(S\) satisfies
\[
 \Lambda-\mu
   =\sum_{\substack{a<c\\\epsilon(a)=\epsilon(c)}}m_{ac}(e_a-e_c),
 \qquad m_{ac}\in\mathbb N .
 \tag{HCU.25}
\]
This is the literal relation \(\mathrm{CDom}(\epsilon,\mu,\Lambda)\).
For \(\Lambda=\mathrm{lowK}(\epsilon,\lambda)\), the original
\(\mathrm{lowKSub}\) is the algebraic supremum of \emph{all} such
subspaces \(S\).  It is not initially assumed finite dimensional or
irreducible.

% BEGIN HCU8
\paragraph{HCU8: smooth vectors, derived actions and a smooth compact-type seed.}
\label{claim:HCU8}
Let \(\sigma\) be a strongly continuous unitary representation of \(U(H)\)
on a complete complex Hilbert space.  Its source smooth vectors form a
complex vector subspace, are preserved by the group and by all
complexified derived operators, and satisfy on that subspace
\[
 \begin{split}
 d\sigma^{\mathbb C}([Z,W])&=
     [d\sigma^{\mathbb C}(Z),d\sigma^{\mathbb C}(W)],\\
 \sigma(g)d\sigma^{\mathbb C}(Z)v
   &=d\sigma^{\mathbb C}(gZg^{-1})\sigma(g)v,\\
 \langle d\sigma^{\mathbb C}(Z)v,w\rangle
   &=-\langle v,d\sigma^{\mathbb C}(c_H(Z))w\rangle .
 \end{split}
 \tag{HCU.26}
\]
Every finite-dimensional irreducible subspace for a compact subgroup
\(K\) has an isomorphic \(K\)-subspace consisting of smooth vectors.
For \(K=K_\epsilon\), this preserves its exact
\(\mathrm{IsKTypeOf}\) weight.  In particular every original DSModel
has a smooth compact-type seed of its specified lowest weight.

\emph{Proof.}
The product-exponential charts established in HCU5 apply equally to
Hilbert-valued functions.  A smooth orbit map remains smooth after
addition, scalar multiplication and translation.  If \(v\) is smooth,
the function \((g,t)\mapsto\sigma(g e^{tX})v\) is smooth.
Differentiating in \(t\) at zero gives the smooth orbit map of
\(d\sigma(X)v\).  This proves stability under differentiation without
an assumption that an unbounded derived operator is continuous on the
whole Hilbert space.  The ordinary chain rule in these charts proves
real linearity in \(X\) and complex linearity in \(v\).  Complexification
gives linearity in every \(Z\).

Differentiate
\(g e^{tX}g^{-1}=e^{t gXg^{-1}}\) for covariance.
The mixed derivatives of
\(e^{sX}e^{tY}e^{-sX}e^{-tY}=1+st[X,Y]+o(st)\), or the
corresponding right-invariant vector fields in a chart, give the
commutator identity.  All derivatives here act on the same smooth
orbit map, so their interchange is legitimate.
Differentiating
\(\langle\sigma(e^{tX})v,\sigma(e^{tX})w\rangle=\langle v,w\rangle\)
gives skew adjointness for real \(X\); decomposing \(Z=X+iY\)
and remembering conjugate linearity in the first argument gives the
last formula in (HCU.26).

Here is the smoothing argument with its exact compact type preserved.
Let \(S\) be the given finite-dimensional irreducible \(K\)-subspace
and \(v\in S\setminus\{0\}\).
The real part of \(\langle v,\sigma(g)v\rangle\) is positive on a
neighborhood of \(1\).  Compactness of \(K\) gives a smaller neighborhood
all of whose \(K\)-conjugates lie in it: continuity of conjugation on
\(K\times\{1\}\) and a finite cover of \(K\) prove this assertion.
Choose a nonnegative smooth bump supported in that smaller neighborhood,
positive at \(1\), and average it under compact conjugation to obtain
\(a\).  The average is smooth, compactly supported, nonnegative and
\(K\)-conjugation invariant, and its support remains in the positivity
neighborhood.  These bumps exist directly in a finite-dimensional
real chart by the usual function
\(e^{-1/(1-\|x\|^2)}\) inside a ball, extended by zero.

Conjugation by \(K\) preserves left Haar measure on \(U(H)\).
Indeed its scaling factor is a continuous homomorphism
\(K\to\mathbb R_{>0}\), whose compact image must be \(\{1\}\);
taking logarithms proves the last assertion.
Define
\[
 A v'=\int_{U(H)}a(g)\sigma(g)v'\,dg\quad(v'\in S).
 \tag{HCU.27}
\]
Strong continuity and compact support give its Hilbert-valued integral.
The real part of \(\langle v,Av\rangle\) is positive, so \(A\ne0\).
Changing variables by compact conjugation shows
\(A\sigma(k)=\sigma(k)A\).  Its kernel in \(S\) is invariant;
irreducibility therefore makes \(A\) injective.  Its image is finite
dimensional and isomorphic to \(S\) as a compact representation, so all
parts of \(\mathrm{IsKTypeOf}\), including (HCU.25), are unchanged.

Finally
\(\sigma(x)Av'=\int a(x^{-1}g)\sigma(g)v'\,dg\).
On a compact neighborhood of any \(x\), this integrand and all its
coordinate derivatives have a common compact support in \(g\);
their norms are bounded by a scalar smooth-bump derivative times
\(\|v'\|\).  Dominated differentiation of the Bochner integral gives
a smooth orbit map for \(Av'\).  The argument is unchanged after
conjugating the matrix group by the original \(M.\mathrm{fr}\).
Since the source \(M.\mathrm{sm}\) is the span of smooth vectors,
the already proved linear-subspace property identifies it with the
smooth vectors themselves.  This proves the complete seed contract
\texttt{exists\_lowK\_le\_sm}.
% END HCU8

% BEGIN HCU9
\paragraph{HCU9: the finite compact-weight facts actually needed.}
\label{claim:HCU9}
The following facts hold for a unitary continuous representation of
\(K_\epsilon\).
\begin{enumerate}
\item Every nonzero finite-dimensional invariant subspace contains an
\(\mathrm{IsKTypeOf}(\epsilon,\sigma,\mu,S')\).
Every such highest weight \(\mu\) is integral and weakly decreasing
along each increasing-index sign block.
\item If \(L\) is the supremum of all compact subspaces of fixed type
\(\Lambda\), then every vector of \(L\) belongs to a finite-dimensional
compact-stable subspace of \(L\).
If \(y\) has torus weight \(\nu\) and
\(\langle w,y\rangle\ne0\) for some \(w\in L\), then
\(\mathrm{CDom}(\epsilon,\nu,\Lambda)\).
\item If a smooth \(v\in L\) has weight \(\Lambda\), then
\(d\sigma^{\mathbb C}(E_{ac})v=0\) for
\(a<c,\epsilon(a)=\epsilon(c)\), and
\(d\sigma^{\mathbb C}(E_{aa})v=\Lambda_a v\).
\end{enumerate}

\emph{Proof.}
A nonzero invariant subspace of smallest positive dimension is
irreducible.  Commuting unitary operators of the diagonal torus on it
are simultaneously diagonalizable: diagonalize one nonscalar operator,
restrict the remaining commuting operators to its eigenspaces, and
induct on dimension.  The resulting simultaneous characters are
products \(e^{i\sum\mu_a s_a}\), with \(\mu_a\in\mathbb Z\).
One elementary verification is to smooth the finite representation by
the averaging argument of HCU8, now on the compact group itself; on an
irreducible finite summand its nonzero intertwiner is invertible.
Each character is therefore smooth, its differential is a real
linear form in the angles, and the ordinary differential equation
along a line gives the exponential formula.
Period \(2\pi\) in each coordinate makes every \(\mu_a\) integral.

Choose a weight maximizing \(\sum_a(-a)\mu_a\) among the finite weight
set, and a nonzero vector of that weight.  Every matrix unit \(E_{ac}\)
with \(a<c\) in the same sign block raises that quantity; hence its
complexified derivative kills this vector.
The span of words in the opposite, lower matrix units applied to it
is invariant under the whole complexified compact Lie algebra.
Here is the algebraic check: move an upper matrix unit to the right
through a word using
\([E_{ab},E_{cd}]=\delta_{bc}E_{ad}-\delta_{da}E_{cb}\);
each commutator shortens the word, and induction reorders it into
lower units, diagonal units and upper units.  Diagonal units act by
the current weight, and upper units at the far right kill the vector.
Thus all resulting weights are the chosen weight minus nonnegative
integer sums of compact positive roots.
The compact group is a product of unitary groups and is connected;
each of its elements is an exponential of a skew-Hermitian block
matrix.  A finite-dimensional Lie-stable subspace is stable under
these exponentials by the matrix exponential series.  Irreducibility
therefore makes this lower-word span the whole summand.
This proves precisely (HCU.25) and all fields of
\(\mathrm{IsKTypeOf}\), without an appeal to a classification theorem.

For any \(i<j\) in a sign block, permutation of these coordinates is
in \(K_\epsilon\).  It carries a vector of weight \(\mu\) to one of
weight \(\mu\circ(ij)\).
Apply (HCU.25) and pair its positive-root sum with the strictly
decreasing vector \((-a)_a\).  The resulting inequality is
\((\mu_i-\mu_j)(j-i)\ge0\), proving dominance.

A vector in an algebraic supremum is a finite sum of vectors from the
defining subspaces; their finite sum is the required finite invariant
space.  The orthogonal projection onto any one defining compact
subspace commutes with the compact action, since its orthogonal
complement is invariant too.  If \(\nu\) fails (HCU.25), the projection
of \(y\) onto every such subspace is zero, by its weight condition.
Then \(y\) is orthogonal to their algebraic sum \(L\), which proves
assertion 2 by contraposition.

For assertion 3, choose a finite compact-stable subspace of \(L\)
containing \(v\).  Differentiation along compact real directions
preserves it because it is closed.  Their complex span contains
\(E_{ac}\) within each sign block.  Covariance in (HCU.26) shows that
\(d\sigma^{\mathbb C}(E_{ac})v\) has weight
\(\Lambda+e_a-e_c\).  This weight cannot satisfy (HCU.25): pairing
\(-e_a+e_c\) with \((-j)_j\) is negative.  Assertion 2, applied to
the derivative vector itself, makes it zero.  Differentiation of the
diagonal torus eigencharacter proves the diagonal formula.
% END HCU9

% BEGIN HCU10
\paragraph{HCU10: the precise lowest-norm inequality.}
\label{claim:HCU10}
Suppose \(\epsilon(a)\) is true exactly when \(a<p\), with
\(0<p<n\), and \(\lambda\) is strictly decreasing.
Write the literal source weights and norm as
\[
 \Lambda=\lambda+\rho_N-\rho_C,\quad
 V(\eta)=\|\eta+2\rho_C\|^2,\quad
 \rho_N(a)+\rho_C(a)=\frac{n-1}{2}-a.
 \tag{HCU.28}
\]
The last equality follows by counting all indices before and after \(a\).
Assume \(\Lambda,\mu\) are dominant on each sign block and
\[
 \mathrm{CDom}(\epsilon,\mu+e_a-e_c,\Lambda),
 \qquad a<p\le c .
 \tag{HCU.29}
\]
Then \(V(\mu)<V(\Lambda)\).

\emph{Proof.}
Put \(L=\Lambda+2\rho_C\), \(M=\mu+2\rho_C\), and
\(\gamma=\Lambda-\mu-e_a+e_c\), a nonnegative sum of compact positive
roots.  The vectors \(\rho_C,L,M\) are weakly decreasing within each
block, by their definitions and dominance.
For \(a_0=p-1,c_0=p\), (HCU.28) gives
\[
 L_{a_0}-L_{c_0}
       =\lambda_{a_0}-\lambda_{c_0}+1>1.
 \tag{HCU.30}
\]
Every compact positive root has nonpositive \(a_0\)-coordinate and
nonnegative \(c_0\)-coordinate.  Therefore
\(\gamma_{a_0}-\gamma_{c_0}\le0\).
Since \(L-M=e_a-e_c+\gamma\), it follows that
\(M_{a_0}-M_{c_0}\ge L_{a_0}-L_{c_0}-2\).
Pairing the root expansion with the block-decreasing vector \(L+M\)
now gives
\[
 \begin{split}
 \|L\|^2-\|M\|^2
 &= (L_a+M_a)-(L_c+M_c)+\langle\gamma,L+M\rangle\\
 &\ge (L_{a_0}+M_{a_0})-(L_{c_0}+M_{c_0})\\
 &\ge 2(L_{a_0}-L_{c_0})-2>0.
 \end{split}
 \tag{HCU.31}
\]
This is the exact \texttt{Lowering.vNormSq\_lt} argument needed below.
For an empty sign block there is no pair \(a<p\le c\); the associated
annihilation statement is vacuous.
% END HCU10

% BEGIN HCU11
\paragraph{HCU11: lowering on the full lowest subspace.}
\label{claim:HCU11}
Suppose \(\mathrm{IsDS}(\epsilon,\lambda,\sigma)\),
\(\mathrm{IsHCParam}(\lambda)\), and the ordered sign pattern above.
Let
\(L=\mathrm{lowKSub}(\epsilon,\lambda,\sigma)\).
For every smooth \(v\in L\),
\[
 d\sigma^{\mathbb C}(E_{ca})v=0\quad(a<p\le c).
 \tag{HCU.32}
\]
This proves the original \texttt{lowK\_killed}, including its use on
the whole supremum \(L\).

\emph{Proof.}
The source \(\mathrm{IsDS}\) supplies a complete unitary representation,
a compact type of weight \(\Lambda=\mathrm{lowK}(\epsilon,\lambda)\),
and the minimality condition
\[
 V(\Lambda)\le V(\mu)
 \quad\hbox{for every actual compact type of weight }\mu.
 \tag{HCU.33}
\]
Choose, by HCU9, a finite compact-stable space \(W\subset L\) containing
\(v\), and replace it by its intersection with the smooth vectors.
This intersection is a vector subspace by HCU8, is finite dimensional,
is compact-stable, and still contains \(v\).
Let \(S\) be the span of
\(d\sigma^{\mathbb C}(Z)w\), where \(w\in W\) and \(Z\) is a linear
combination of \(E_{ca}\) with \(a<p\le c\).
It is finite dimensional, being the image of the finite-dimensional
space of such \(Z\)'s tensor \(W\).
It consists of smooth vectors and is compact-stable by covariance,
since compact block conjugation preserves these off-diagonal units.

If one expression in (HCU.32) is nonzero, then \(S\ne0\).
HCU9 gives an irreducible compact subspace \(S'\subset S\) of highest
weight \(\mu\), with a nonzero weight vector \(u\).
Some generator has
\(\langle d\sigma^{\mathbb C}(E_{ca})w,u\rangle\ne0\); otherwise
\(u\), which belongs to their span, would be orthogonal to itself.
Since \(c_{J_\epsilon}(E_{ca})=E_{ac}\), the last identity in
(HCU.26) makes
\[
 \langle w,d\sigma^{\mathbb C}(E_{ac})u\rangle\ne0.
 \tag{HCU.34}
\]
(The adjoint identity has a minus sign, which does not affect
nonvanishing.)  Covariance under the diagonal torus gives the latter
vector weight \(\mu+e_a-e_c\).
Because \(w\in L\), HCU9 implies (HCU.29).
Both \(\mu\) and \(\Lambda\) are compact dominant by HCU9.
HCU10 contradicts (HCU.33).  Thus \(S=0\), proving (HCU.32).
No prior finite-dimensionality or multiplicity-one assertion for \(L\)
has entered the proof.
% END HCU11

% BEGIN HCU12
\paragraph{HCU12: original frames, extremal vectors and Borel equations.}
\label{claim:HCU12}
Let \(M\) be any original DSModel with \(\mathrm{HolShape}(M)\).
At each real place put \(q=n-p\).
Let \(r=\mathrm{revNeg}(\epsilon)\), the permutation reversing the
negative sign block and fixing the positive block, and let
\[
 a_q(j)=
 \begin{cases}n-1-j&j<q,\\j-q&j\ge q.\end{cases}
 \quad
 x_0=M.\mathrm{fr}^{-1}P_{a_q},\quad
 w_j=\Lambda_{\,r(a_q(j))}.
 \tag{HCU.35}
\]
These are the source's \(\mathrm{arr}\), \(\mathrm{x0Mat}\) and
\(\mathrm{WtIs}\) formulas.  If \(v\) is a smooth vector in the full
lowest subspace and has the extremal torus weight
\(\Lambda\circ r\), then the original \(M.\mathrm{rep}\) satisfies
\[
 d(M.\mathrm{rep})^{\mathbb C}(x_0E_{ac}x_0^{-1})v
       =\delta_{ac}w_av\qquad(a\le c).
 \tag{HCU.36}
\]
In particular every supplied source \(\mathrm{IsExtVec}\) satisfies
the full \texttt{borel\_extVec} conclusion.

\emph{Proof.}
The stored identity
\(M.\mathrm{fr}^*J_\epsilon M.\mathrm{fr}=H\) gives
\(x_0^*Hx_0=\operatorname{diag}(-1_q,1_p)\), exactly (HCU.1).
Conjugating the representation by \(M.\mathrm{fr}\) preserves
unitarity, irreducibility and smoothness on their original Hilbert
space.  Direct multiplication gives
\[
 x_0E_{ac}x_0^{-1}
    =M.\mathrm{fr}^{-1}E_{a_q(a),a_q(c)}M.\mathrm{fr}.
 \tag{HCU.37}
\]
The same identity for derived operators follows by differentiating
\(F e^{tX}F^{-1}=e^{tFXF^{-1}}\); its complexification respects
the original Hermitian involutions.

For a diagonal matrix unit, HCU9 gives the required extremal weight.
For a strict upper unit in the two negative-first coordinates of
(HCU.35), there are three cases.
If both indices lie in the negative block, their images under
\(r\circ a_q\) are \(p+a,p+c\), in increasing order.
If both lie in the positive block, these images are \(a-q,c-q\),
again in increasing order.  The permutation \(r\) is compact;
conjugating \(v\) by its permutation matrix turns it into a vector
of weight \(\Lambda\).  HCU9's compact raising annihilation applies
in both cases.  In the remaining case \(a<q\le c\), the matrix unit
in (HCU.37) has its row in the negative block and column in the
positive block, so HCU11 applies.
This proves (HCU.36) with the actual order and signs.

For use in constructing an extremal vector, HCU8 supplies a smooth
compact subspace of highest weight \(\Lambda\); applying the compact
permutation \(r\) to a normalized highest vector gives a unit vector
of weight \(\Lambda\circ r\) in that same subspace.
Its torus weights are integral by HCU9, so the \(w\)'s in (HCU.35)
are integers.  This produces the required weight tuple before
any appeal to lowest-type multiplicity one.
% END HCU12

% BEGIN HCU13
\paragraph{HCU13: full LowSm and LowIrr, with multiplicity one proved.}
\label{claim:HCU13}
For every original DSModel \(M\) with \(\mathrm{HolShape}(M)\),
\[
 \mathrm{LowSm}(M)\ \wedge\ \mathrm{LowIrr}(M).
 \tag{HCU.38}
\]
Thus the entire original
\texttt{HolLowOnce.lowSm\_and\_lowIrr} producer is reconstructed.

\emph{Proof.}
Work at one real place.  HCU8 supplies a smooth irreducible compact
subspace \(S_0\) of weight \(\Lambda\).
Choose its unit extremal vector \(v\) by HCU12.
Let \(T\) be any other subspace satisfying
\(\mathrm{IsKTypeOf}(\epsilon,\sigma,\Lambda,T)\).
Its definition supplies a nonzero highest vector, and the compact
permutation \(r\), followed by normalization, gives a unit extremal
\(v'\in T\) of weight \(\Lambda\circ r\).
This \(v'\) is not assumed smooth for the full noncompact group.

Set \(l=v'-\langle v,v'\rangle v\).
It has the same torus weight, and \(\langle l,v\rangle=0\), by
conjugate symmetry and \(\|v\|=1\).
The representation on the original Hermitian coordinates is unitary
and irreducible by its stored DSModel identities.
HCU12 gives the Borel equations for the smooth vector \(v\), and
(HCU.35) identifies the torus character of both \(v\) and \(v'\)
with the exact \(\chi_w\).
HCU7 therefore gives \(l=0\), so \(v'\in S_0\).

The intersection \(T\cap S_0\) is nonzero and compact-stable.
Irreducibility first gives \(T\subset S_0\), and then \(T=S_0\).
Taking the supremum over \emph{all} defining subspaces proves
\(\mathrm{lowKSub}=S_0\).
It is consequently contained in the smooth vectors and is itself
an \(\mathrm{IsKTypeOf}\) of the required weight.  These are exactly
\(\mathrm{LowSm}\) and \(\mathrm{LowIrr}\), place by place.
The equality was proved on the original Hilbert space and with the
original compact action, not by selecting a different DSModel.
% END HCU13

\subsection*{Actual real lifts, holomorphic forms and the Hilbert coefficient}

% BEGIN HCU14
\paragraph{HCU14: the original archimedean lift, with no changed coordinates.}
\label{claim:HCU14}
For the original \(D,\theta,V\) with \(\mathrm{IsInfFam}(\theta)\),
the map
\[
 G_{\infty}(V)\longrightarrow
 \prod_\tau U(\mathrm{gramAt}(D,V,\tau)),\qquad
 g\longmapsto(\mathrm{archM}(V,\theta_\tau,g))_\tau
 \tag{HCU.39}
\]
is a continuous group isomorphism with continuous inverse \(L\).
The inverse of a tuple supported at \(\tau\) is the original
\(\mathrm{atTau}(\theta,V,\tau,g_\tau)\).
Every adelic \(g\) decomposes as \(g=au\) with
\(a\in G_f(V)\), \(u\in G_\infty(V)\), and its archimedean
evaluations belong to the stated unitary groups.
The finite and infinite adelic subgroups commute, and each original
\(\mathrm{gramAt}\) is Hermitian and nonsingular.
These are all the literal \texttt{exists\_real\_lift} fields.

\emph{Proof.}
Write \(E/F\) for the original CM quadratic extension.
The stored \(\mathrm{tauE}\) indices give exactly one member of each
conjugate pair of embeddings of \(E\).  More explicitly, the source
maps its index set bijectively to the chosen set \(\Psi\) of embeddings,
and the defining CM-type condition is
\(\sigma\in\Psi\) if and only if \(\overline{\sigma}\notin\Psi\).
Two embeddings of \(E\) over a real embedding of \(F\) are such a
pair.  Thus these indices enumerate the real places of \(F\) once.

Here the infinitary homomorphisms are forced by their stated
\(\mathrm{IsInfEmb}\) conditions, not chosen afresh.
The diagonal copy of \(F\) is dense in \(F_\infty=\prod_{v\mid\infty}F_v\).
One direct proof chooses a rational basis of \(F\); its real embedding
matrix is invertible by separability, as is seen from the Vandermonde
matrix of a primitive element and change of basis.
The image of the dense set \(\mathbb Q^{[F:\mathbb Q]}\) under this
real linear isomorphism is the required diagonal copy.
A homomorphism satisfying \(\mathrm{IsInfEmb}(\tau,\theta_\tau)\)
therefore agrees on \(F_\infty\) with evaluation at the real place
under \(\tau\), by continuity and its specified values on \(F\).
For a pure tensor it must consequently satisfy
\[
 \theta_\tau(a\otimes e)=a_{v(\tau)}\,\tau(e).
 \tag{HCU.40}
\]

To verify that these evaluations are a bijection, take
\(\delta=y-\overline y\ne0\) in \(E\).
The vectors \(1,\delta\) form an \(F\)-basis, since conjugation fixes
the first and negates the second.  At every real place,
\(\tau(\delta)\) is nonzero purely imaginary.
Thus the two real coordinates of \(a+b\delta\) give all of
\(\mathbb C\) at that place, uniquely and continuously.
Place by place, (HCU.40) is the continuous ring isomorphism
\[
 F_\infty\otimes_F E\simeq\prod_\tau\mathbb C
 \tag{HCU.41}
\]
that takes the original CM involution to complex conjugation.
The formula is the given \(\theta\), not just an isomorphic alternative.

Apply this isomorphism to each matrix entry.  The original Gram matrix
evaluates to \(\mathrm{gramAt}\); the unitary equation and inverse
matrices are preserved in both directions.
The original \(HSp\) identity that its Gram matrix equals its
CM-conjugate transpose evaluates to ordinary Hermitian symmetry.
Its nonzero determinant stays nonzero under the field embedding;
thus each evaluated Gram matrix is nonsingular.
In particular evaluating any adelic unitary \(g\), without assuming
that its finite component is the identity, gives an element of
each stated \(U(\mathrm{gramAt})\).
The adelic ring is the product of its finite and infinite parts.
Gluing the tuple in (HCU.41) to the identity finite matrix therefore
defines the unique element \(L(u)\in G_\infty(V)\).
It is continuous and multiplicative, and both inverse identities in
(HCU.39) follow entry by entry.
The stored \(\mathrm{atTau}\) has exactly the same finite identity
and single-place evaluations, so uniqueness identifies it with
\(L(\mathrm{single}_\tau g_\tau)\).
Finally matrices supported in the two different product factors
commute.  For an arbitrary adelic unitary \(g\), retain its finite
component and replace its infinite component by the identity to
obtain \(a\); perform the opposite replacement to obtain \(u\).
Both satisfy the original Gram equation componentwise, and \(g=au\).
This supplies the asserted finite--infinite decomposition.
These verifications retain the original \(F,E,\theta\),
Hermitian equation and all evaluation maps.
% END HCU14

% BEGIN HCU15
\paragraph{HCU15: the actual extremal model coefficient is holomorphic.}
\label{claim:HCU15}
Let \(M\) be an original DSModel, and assume its literal
\(\mathrm{HolShape}\), integer-weight \(\mathrm{WtIs}(M,w)\)
and \(\mathrm{IsExtVec}(M,v)\).
Then
\[
 B(u)=\prod_\tau
 \langle v_\tau,M.\mathrm{rep}_\tau(u_\tau)v_\tau\rangle
 \tag{HCU.42}
\]
has \(\mathrm{HolFn}(\mathrm{gramAt},q,\mathrm{x0Mat},w)\),
satisfies the left torus law of HCU4, and \(B(1)=1\).
This proves the whole \texttt{holFn\_coef\_extVec} output used
by the original L2 argument.

\emph{Proof.}
HCU8 identifies each \(v_\tau\in M.\mathrm{sm}_\tau\) with a genuine
smooth vector.  Differentiation of a coefficient gives
\[
 R_Z^{\mathbb C}
   \langle v_\tau,M.\mathrm{rep}_\tau(u)v_\tau\rangle
 =\langle v_\tau,M.\mathrm{rep}_\tau(u)
                 d(M.\mathrm{rep}_\tau)^{\mathbb C}(Z)v_\tau\rangle.
 \tag{HCU.43}
\]
This follows first for real \(Z\) by bounded linearity of the
fixed group operator and the inner product functional, and then
by the complexification (HCU.16).
HCU12 supplies the exact Borel equations for (HCU.43).
The finite product (HCU.42) thus has the source \(\mathrm{SmFn}\)
and \(\mathrm{KilledFn}\); HCU5 makes it a \(\mathrm{HolFn}\).

The original extremal weight and the permutation in (HCU.35) give
\[
 M.\mathrm{rep}_\tau(t(s)_\tau)v_\tau
 =\left(\prod_a e^{iw_{\tau,a}s_{\tau,a}}\right)v_\tau.
\]
Unitarity, as in (HCU.24), gives the left torus law for each
coefficient and hence their product.  The normalization in
\(\mathrm{IsExtVec}\) is \(\|v_\tau\|=1\), so \(B(1)=1\).
Every factor in (HCU.42) is the original model coefficient.
% END HCU15

% BEGIN HCU16
\paragraph{HCU16: original finite-volume representatives.}
\label{claim:HCU16}
Fix the original \(\mathrm{FinPair}(V,\mu,\mathcal F)\) and
\(\mathrm{IsReg}(V,\mu,\mathcal F,\rho)\), and put
\(\nu_Q=(\mu|_{\mathcal F})_*\mathrm{autMk}\).
A continuous bounded rational-left-invariant function \(\phi\)
has an \(L^2(\nu_Q)\) representative \(F\) satisfying the source
\(\mathrm{Repr}(F,\phi)\).
Representatives obey linearity and right translation, and
\[
 \langle F,F'\rangle
   =\int_{\mathcal F}\overline{\phi(x)}\phi'(x)\,d\mu(x)
 \tag{HCU.44}
\]
whenever \(\mathrm{Repr}(F,\phi)\) and
\(\mathrm{Repr}(F',\phi')\).
The integrand is integrable.  The \(L^2\) representative is unique;
a continuous representative is also unique pointwise.
In particular
\[
 \phi(xk)=c\phi(x)\ (\forall x)
 \quad\Longrightarrow\quad \rho(k)F=cF .
 \tag{HCU.45}
\]

\emph{Proof.}
The invariant \(\phi\) descends to a measurable function on the
quotient with its quotient measurable structure.  The pushed-forward
measure is finite, since \(\mu(\mathcal F)<\infty\).
Boundedness makes the descended function square integrable.
Its chosen \(L^2\) representative equals it almost everywhere on the
quotient, hence almost everywhere on \(\mathcal F\) after pullback.
The exceptional set in the full group is rational-left-invariant.
Countably many rational translates of the fundamental domain cover
the group up to a null set; the rational group is countable since
it is a set of matrices over the number field.
Left invariance of \(\mu\) therefore makes that exceptional set
globally null.  This is exactly the source's globally quantified
\(\mathrm{Repr}\), not merely an equality on the fundamental domain.

Conversely equality of two representatives globally almost everywhere
implies their equality for \(\mu|_{\mathcal F}\) and then for \(\nu_Q\);
this proves \(L^2\) uniqueness.
If two continuous functions agree almost everywhere for Haar measure,
their difference is zero everywhere: a nonzero value would give a
nonempty open set on which its modulus is bounded below, and Haar
measure is positive on every nonempty open set.
The usual almost-everywhere identities for \(L^2\) addition and scalar
multiplication lift globally by the same invariant-null-set argument.
The \(\mathrm{IsReg.val}\) identity and the right invariance in
\(\mathrm{FinPair}\) give
\(\mathrm{Repr}(\rho(k)F,\phi(\,\cdot\,k))\).
Finally Cauchy--Schwarz gives integrability of the product of the two
\(L^2\) functions.  The definition of the pushforward measure and
the representative identities transform its inner product into
(HCU.44).  These facts prove linearity, translation and (HCU.45),
including zero vectors and zero scalar eigenvalues where applicable.
No existence of a finite-volume target quotient is asserted here:
the same \(\mathrm{FinPair}\) remains an explicit input.
% END HCU16

% BEGIN HCU17
\paragraph{HCU17: the original extremal coefficient identity.}
\label{claim:HCU17}
On the original \(D,\theta,V,M,w,v\), assume
\(\mathrm{IsInfFam}\), \(\mathrm{HolShape}\), \(\mathrm{WtIs}\),
\(\mathrm{IsExtVec}\), the same \(\mathrm{FinPair}\) and
\(\mathrm{IsReg}\).
Assume also \(\mathrm{HolForm}(\theta,M,w,J,\chi,\phi)\)
and \(\mathrm{Repr}(F,\phi)\).
Then for every \(g\in G_\infty(V)\),
\[
 \langle F,\rho(g)F\rangle
 =\|F\|^2\prod_\tau
  \langle v_\tau,M.\mathrm{rep}_\tau
          (\mathrm{archM}(V,\theta_\tau,g))v_\tau\rangle .
 \tag{HCU.46}
\]
This is the entire literal
\texttt{inner\_reg\_eq\_coef\_extremal} statement.

\emph{Proof.}
Use the original \(L\) of HCU14.
The \(\mathrm{HolAd}\) field of \(\mathrm{HolForm}\) says that, for
each fixed \(x\),
\(u\mapsto\phi(xL(u))\) has the exact \(\mathrm{HolFn}\) type.
If \(|\phi|\le C\), then
\[
 (x,u)\longmapsto\overline{\phi(x)}\phi(xL(u))
 \tag{HCU.47}
\]
is bounded in modulus by \(\max(C,0)^2\).
It is measurable in \(x\), since \(L(u)\) is fixed when this
measurability is required and \(\phi\) is continuous.
Apply HCU6 with the finite measure \(\mu|_{\mathcal F}\).
By HCU16 the resulting holomorphic-type function is
\(A(u)=\langle F,\rho(L(u))F\rangle\).

The right upper-triangular law (HCU.4), evaluated at the diagonal
torus, gives directly
\[
 \phi(xL(t(s)))=\chi_w(s)\phi(x).
\]
Thus \(\rho(L(t(s)))F=\chi_w(s)F\) by (HCU.45).
Unitarity of \(\rho\) gives
\(A(t(s)u)=\chi_w(s)A(u)\), with the same character orientation as
(HCU.24), and \(A(1)=\|F\|^2\).
HCU15 gives the function \(B\) in (HCU.42), with the identical
holomorphic type and left torus law, and \(B(1)=1\).
HCU7 therefore yields \(A=\|F\|^2B\).
Use both inverse identities of HCU14 to replace \(L(u)\) by the
given \(g\in G_\infty(V)\); this gives (HCU.46).
The conclusion is valid when \(F=0\).
The boundedness of \(\phi\) and finite target measure were used as
explicit inputs, not inferred from source compactness or from the
name of a rapid-decay theorem.
% END HCU17

% BEGIN HCU18
\paragraph{HCU18: the exact Hilbert multilinear realization.}
\label{claim:HCU18}
Let \(G_\tau=U(H_\tau)\) act unitarily and irreducibly on complete
Hilbert spaces \(\mathcal H_\tau\), with unit vectors \(v_\tau\).
Let \(\rho\) be a unitary representation of
\(G=\prod_\tau G_\tau\) on a complete Hilbert space \(\mathcal K\).
Suppose \(k\in\mathcal K\) satisfies
\[
 \langle k,\rho(u)k\rangle
  =\|k\|^2\prod_\tau\langle v_\tau,\sigma_\tau(u_\tau)v_\tau\rangle .
 \tag{HCU.48}
\]
Then there is a continuous multilinear map
\(\Phi:\prod_\tau\mathcal H_\tau\to\mathcal K\) with
\[
 \begin{aligned}
 \Phi(v)&=k,\\
 \Phi(y_1,\ldots,\sigma_\tau(g)y_\tau,\ldots)
      &=\rho(\mathrm{single}_\tau g)\Phi(y),\\
 \|\Phi(y)\|&=\|k\|\prod_\tau\|y_\tau\|.
 \end{aligned}
 \tag{HCU.49}
\]
and whose image lies in the closed span of \(\rho(G)k\).

This proves the full source
\texttt{exists\_equivariant\_}\allowbreak
\texttt{multilinear\_of\_}\allowbreak\texttt{matrixCoeff\_eq}
contract.

For an adelic \(\rho\), HCU14 transports this construction to the
original one-place equivariance with \(\mathrm{atTau}\).
Every finite-place eigenlaw \(\rho(a)k=c k\), \(a\in G_f(V)\),
then holds for every \(\Phi(y)\).  Hence, with the inputs of HCU17,
the whole original \texttt{BackForm.exists\_map} conclusion holds
for the same \(L^2\) space and same initial representative \(F\).

\emph{Proof.}
Irreducibility makes the span
\(E_\tau=\operatorname{span}\{\sigma_\tau(g)v_\tau:g\in G_\tau\}\)
dense in \(\mathcal H_\tau\), because its closure is nonzero and
invariant.  The inverse in each group and unitarity transform
(HCU.48) into the orbit Gram identity
\[
 \langle \rho(a)k,\rho(b)k\rangle
   =\|k\|^2\prod_\tau
       \langle\sigma_\tau(a_\tau)v_\tau,
              \sigma_\tau(b_\tau)v_\tau\rangle .
 \tag{HCU.50}
\]
For finite expressions
\(y_\tau=\sum_j c_{\tau,j}\sigma_\tau(g_{\tau,j})v_\tau\),
define
\[
 \Phi_0(y)=
  \sum_{(j_\tau)_\tau}
       \left(\prod_\tau c_{\tau,j_\tau}\right)
       \rho((g_{\tau,j_\tau})_\tau)k .
 \tag{HCU.51}
\]
Expanding (HCU.50) gives
\[
 \langle\Phi_0(y),\Phi_0(z)\rangle
       =\|k\|^2\prod_\tau\langle y_\tau,z_\tau\rangle .
 \tag{HCU.52}
\]
If two expressions represent the same tuple, expansion of the squared
norm of the difference using (HCU.52) gives zero.  Thus (HCU.51)
is well defined and multilinear, and its norm is the one in
(HCU.49).  Approximate an arbitrary tuple by tuples from the dense
\(E_\tau\)'s.  A telescoping multilinear difference, bounded by the
product norm formula, proves that the images are Cauchy.
Completeness gives a limit independent of all approximations and
yields the continuous multilinear \(\Phi\).
All stated identities pass from the finite sums to these limits.
This includes the image in the closed orbit span.  If \(k=0\), the
construction gives the zero map; if the index set is empty it gives
the constant value \(k\), the correct zero-variable multilinear map.

For the adelic conclusion, compose the actual \(\rho\) with the
real lift \(L\) and use (HCU.46).
The original model representations are unitary, irreducible and
complete by the DSModel fields and their fixed frame conjugations.
HCU14 identifies the one-place action with \(\mathrm{atTau}\).
Since a finite-place \(a\) commutes with every \(L(u)\), its
\(c\)-eigenlaw holds on every real orbit vector \(\rho(L(u))k\)
and then on their closed span by boundedness of \(\rho(a)-c\).
HCU16 transfers each pointwise finite eigenlaw of the original
\(\phi\) to its \(L^2\) representative \(F\).
These are precisely the four outputs of
\texttt{BackForm.exists\_map}, including its original finite
characters and its larger full-adelic closed-orbit containment.
This is a Hilbert-valued map; conversion to pointwise smooth
automorphic functions is a separate remaining step.
% END HCU18

% BEGIN HCU19
\paragraph{HCU19: the original lowest-type cusp packet is holomorphic at its extremal.}
\label{claim:HCU19}
Let \(M,\theta,w\) satisfy the original
\(\mathrm{HolShape}\), \(\mathrm{IsInfFam}\) and \(\mathrm{WtIs}\).
Let \(U\) be an original lowest-type multilinear map whose values
are smooth, whose compact equivariance is the source
\(\mathrm{KEquiv}\), and which is killed by every original
\(\mathfrak p^-\) direction.
If \(e\) is a tuple with
\(\mathrm{IsExtVec}(M,(e_\tau)_\tau)\), then
\[
 \mathrm{HolAd}(\theta,M,w,U(e)).
 \tag{HCU.53}
\]
Thus the complete \texttt{holAd\_of\_kEquiv} consumer is reconstructed
with LowSm and LowIrr supplied by HCU13.

\emph{Proof.}
The lowest compact subspaces are finite dimensional by HCU13.
For any fixed output point, evaluation of the multilinear map \(U\)
is therefore a continuous multilinear form on them.
Differentiate its compact equivariance along a real compact
one-parameter group.  Ordinary finite-dimensional differentiation
gives its derivative by replacing the corresponding input by the
derived compact action.  Complexify within the compact blocks.
No continuity of \(U\) on the full discrete-series Hilbert spaces
is being assumed.

Fix \(g\) in the target adelic group and consider
\(f_g(u)=U(e)(gL(u))\), with the original \(L\) of HCU14.
It satisfies the source \(\mathrm{SmFn}\).
For the diagonal and same-sign upper units in the \(x_0\) coordinates,
the preceding differentiation and the compact part of HCU12 give
exactly \(w_a f_g\) and zero, respectively.
The remaining upper units are precisely the original
\(\mathfrak p^-\) directions after (HCU.37), and vanish by the
assumed lowering equations.
Thus \(f_g\) satisfies (HCU.17).  HCU5 gives its literal
\(\mathrm{HolFn}\).
Since \(L\) inverts the original archimedean evaluation, this is
exactly the \(\mathrm{HolAd}\) requirement for every \(g\).
% END HCU19

% BEGIN HCU20
\paragraph{HCU20: installation on the same original base model.}
\label{claim:HCU20}
Let \(R:\mathrm{RealDat}(D,A,K_s)\) be the original datum and retain
its exact \(\mathrm{models\_order}\) field.  For every original \(i\),
put \(M_i=R.MW(i,\mathrm{sg}_0)\).  Then
\[
 \mathrm{HolShape}(M_i),\qquad
 \exists w_i,\ \mathrm{WtIs}(M_i,w_i),\qquad
 \mathrm{LowSm}(M_i)\ \wedge\ \mathrm{LowIrr}(M_i).
 \tag{HCU.54}
\]
There is also a unit original extremal tuple, with the same integer
weight \(w_i\).  This supplies the entire
\texttt{holShape\_fixed} installation and the full-lowest-type inputs
retained in TAE12--TAE13.  It makes no assertion about an arbitrary
other coefficient embedding.

\emph{Proof.}
The field \(R.hv0\) says that \(A.v_0\) is the infinite place of
\(D.\mathrm{what}_2.\mathrm{embLL}\).
By the definition of \(\mathrm{sgAt}\),
\(\mathrm{sg}_0=\mathrm{sgAt}(\mathrm{embLL},A.v_0)=\mathrm{embLL}\).
The stored identity \(D.\mathrm{embE}_{\phi_0}\) makes
\[
 \mathrm{labE}(\mathrm{embLL},\tau)
   =\phi_0\circ\mathrm{sigmaOf}(D.eTau(\tau))
   =\mathrm{tauE}(D,\tau).
 \tag{HCU.55}
\]
The chosen CM-type embeddings are distinct and contain no conjugate
pair, as checked in HCU14.  Therefore any index selected in
\(\mathrm{lab}(\mathrm{embLL},\rho)\) must equal \(\rho\);
the conjugate alternative in its definition is impossible.
Unfolding \(\mathrm{kapAt}\) and \(\mathrm{parAt}\) now gives
\(\mathrm{parAt}(W,\mathrm{sg}_0,\rho)
=\mathrm{holPar}(p,q,\kappa)\), with the original
\(p,q=\mathrm{sig}_W(\mathrm{tauE}(D,\rho))\),
\(p+q=\mathrm{rank}(W)=n\), and
\(\kappa=\mathrm{wtOf}(D,W,\rho)\in\mathbb Z^{p+q}\).
No relabeling or conjugation of the chosen model has occurred.

The model's stored \(\mathrm{enum}\) already asserts that
\(\lambda=M_i.\mathrm{lam}(\rho)\) is strictly decreasing and
enumerates this signed finite parameter set.
Its entries are consequently among
\[
 \begin{aligned}
  \kappa_c^+-\frac q2+\frac{p-1}{2}-c
      &=\frac{n-1}{2}+(\kappa_c^+-q-c)
        &&(0\le c<p),\\
  \kappa_c^-+\frac p2+\frac{q-1}{2}-c
      &=\frac{n-1}{2}+(\kappa_c^--c)
        &&(0\le c<q).
 \end{aligned}
 \tag{HCU.56}
\]
Thus for each index \(a\) there is an integer \(m_a\) with
\(\lambda_a=(n-1)/2+m_a\).  Together with strict decrease this is
exactly \(\mathrm{IsHCParam}(\lambda)\).
The original \(\mathrm{models\_order}\) states
\(\epsilon_a=\mathbf 1_{a<p}\); since \(p\le n\), the remaining
\(\mathrm{HolShape}\) field follows.

For an explicit integral weight, let
\(C_a^+=\#\{b>a:\epsilon_b=\epsilon_a\}\) and
\(C_a^-=\#\{b<a:\epsilon_b=\epsilon_a\}\).
The elementary partition of indices above and below \(a\) gives
\(\rho_N(a)+\rho_C(a)=(n-1)/2-a\).  Hence
\[
 \Lambda_a=\lambda_a+\rho_N(a)-\rho_C(a)
       =m_a+n-1-a-C_a^++C_a^-\in\mathbb Z .
 \tag{HCU.57}
\]
Take \(w_i(\rho,j)=\Lambda_{r(a_q(j))}\), with the actual
\(\mathrm{revNeg}\) and \(\mathrm{arr}\) of HCU12.
This is exactly \(\mathrm{WtIs}\), including empty signature blocks.
HCU13 now gives both full-lowest-type assertions on this same \(M_i\).
HCU8 and HCU12 provide its unit extremal vector at each real place;
their finite tuple uses the just computed weight.
All conclusions hold for every \(i\) before any choice of the
parameter \(t\).  They use \(R.\mathrm{models\_order}\), which is a
field of the already required \(\mathrm{IsModels}\), and do not
construct a new datum \(R\).
% END HCU20

\subsection*{What this closes, and what still must be produced}

The proved endpoints are the actual local and global coefficient
uniqueness, the complete smooth Borel--holomorphic dictionary,
bounded holomorphic integration, full lowest-type smoothness and
multiplicity one, the original extremal coefficient identity,
and its Hilbert multilinear realization with original finite
eigenlaws.  In particular the full-lowest-type producer retained
in TAE is now supplied under its exact \(\mathrm{HolShape}\)
hypothesis; it is not replaced by an existentially selected
irreducible summand.

For the base original model \(M_i=R.MW(i,\mathrm{sg}_0)\), the required
\(\mathrm{HolShape}\) and integer weight are supplied by its ordered
holomorphic parameter formulas, reconstructed with the
\(\mathrm{sg}_0=D.\mathrm{what}_2.\mathrm{embLL}\) identification
in HCU20.  For any other coefficient embedding, the
same HCU implications apply when \emph{that model's}
\(\mathrm{HolShape}\) is checked; this chapter does not assert that
every embedding has that shape.

The original TAE13 cusp packet, together with HCU13 and HCU19,
therefore reaches a holomorphic extremal function.
If the actual boundedness and finite target quotient are supplied,
HCU16--HCU18 give the literal coefficient and Hilbert-map stages.
They do not manufacture those analytic inputs.
The source's predicate \(\mathrm{Published.CuspRapidDecay}\)
has not been turned into a checked published application.
A finite-volume \(\mathrm{FinPair}\), a regular realization
\(\mathrm{IsReg}\), and the original scalar function are kept
on the same target quotient throughout.

To finish \texttt{exists\_multMap}, one still needs the original
smooth representatives and all their derivative identities in that
Hilbert image, point-evaluation Sobolev bounds and resulting
uniform growth, central finiteness, and preservation of cusp
integrals under closed orbit span.  HCU8 supplies smooth-vector
calculus, but it does not prove those global automorphic
point-evaluation estimates.
The remaining exact sources are:
\begin{itemize}
\item \texttt{SmoothVectorAutomorphicDerivatives};
\item \texttt{SobolevBoundForSmoothVectors};
\item \texttt{GelfandInvariantsActByScalarsHighestWeight};
\item \texttt{CuspidalityPreservedInClosure}.
\end{itemize}
No full \(\mathrm{MMap}\), \(\mathrm{firesAt}\) or
\(\mathrm{exists\_thetaMap}\) conclusion is claimed yet.

Independently, one must still construct, on the same original \(R\),
one \(io_i\) with
\[
 \exists io_i\ \forall t,\quad
 \mathrm{IsIoW}\bigl(R.\theta,M_i,R.cV(i),
       R.\omega P(i,t),R.ptP(i,t),\mathrm{Delta}(i),io_i\bigr).
 \tag{HCU.58}
\]
The two intertwining laws, exact polynomial normalization and the
quantifier order are unchanged.  Only such a simultaneous witness
activates the stored \(\mathrm{OrZero}(\mathrm{PIoP}(i))\).
No argument above supplies this polynomial map.
The same-Bundle weak-data producer, geometric Hilbert reading,
rational adjoints and tests, and original \(\mathrm{RowsHold0_3}\)
remain separate.  The present proofs are pointwise and introduce
no new family threshold.  Formal correctness is assumed as in the
project; no new Lean replay is claimed.





